\section{A condition-free best-of-policy wrapper}
\label{sec:portfolio}

The nonidentifiability result concerns a controller forced to choose one
future policy.  Branching removes that information bottleneck.  Let
$P_1,\ldots,P_K$ be local solvers for the same instance and terminal
certificate.  They may include, for example, exact Gauss--Seidel, fixed
two-rung schedules at several band widths, an alpha-scaled ladder, and the
frontier adaptive ladder.  Each policy receives an independent sparse state
initialized from the same input.

For instance $I$, let $C_j(I)$ be the charged edge or row work policy $P_j$
would need if run alone until the verifier accepts it.  Fix positive service
shares $\nu_j$ with $\sum_j\nu_j=1$.  A fractional work scheduler gives arm
$j$ exactly a $\nu_j$ fraction of aggregate charged work and stops all arms
as soon as one certifies.

\begin{theorem}[Weighted portfolio oracle inequality]
\label{thm:portfolio}
Under fractional charged-work scheduling,
\begin{equation}
 C_{\rm port}(I)
 =\min_{1\leq j\leq K}\frac{C_j(I)}{\nu_j}.
 \label{eq:portfolio-exact}
\end{equation}
In particular, uniform service satisfies
\begin{equation}
 C_{\rm port}(I)=K\min_j C_j(I).
 \label{eq:uniform-portfolio}
\end{equation}

If work is scheduled in indivisible chunks and a fair scheduler maintains
received work $w_j(W)\geq\nu_jW-\Delta$ for every arm after aggregate work
$W$, then
\begin{equation}
 C_{\rm port}(I)
 \leq\min_j\frac{C_j(I)+\Delta}{\nu_j}.
 \label{eq:portfolio-discrete}
\end{equation}
\end{theorem}

\begin{proof}
After aggregate fractional work $W$, arm $j$ has received $\nu_jW$.  It
finishes exactly when $W=C_j(I)/\nu_j$; the first finishing arm proves
\eqref{eq:portfolio-exact}.  Uniform shares give
\eqref{eq:uniform-portfolio}.  Under the discrete invariant, aggregate work
$W=(C_j(I)+\Delta)/\nu_j$ guarantees
$w_j(W)\geq C_j(I)$, so arm $j$ has finished.  Minimize over $j$.
\end{proof}

The fractional model is natural for the manuscript's adjacency-edge work:
an adjacency scan can be paused and resumed after buffering its row update.
With atomic service operations, $\Delta$ can instead be taken as the largest
scheduling chunk.  The wrapper uses $K$ sparse solver states and may explore
the union of their frontiers; that storage and scalar-copy overhead must be
reported even though \eqref{eq:portfolio-exact} counts graph work exactly.
The benchmark service used by the measured artifact exposes only one mutable
state, so literal implementation there would require protocol support for
checkpoints or independent instances.
